% ============================================================
\chapter{Diseño e implementación}
\label{cap:diseno}
% ============================================================

Este capítulo describe el diseño completo del sistema \textit{kubescan} y su
implementación como herramienta de software distribuible. Se presentan, en primer lugar,
los requisitos funcionales y no funcionales que guiaron las decisiones de diseño;
a continuación se detalla la arquitectura de tres capas del sistema, los procesos de
construcción del conjunto de datos y del grafo de clúster, y los algoritmos de
entrenamiento de cada capa; finalmente se describe la interfaz de línea de comandos
resultante. La implementación completa se encuentra disponible en el repositorio
indicado en el Apéndice~\ref{ap:repositorio}.

% -----------------------------------------------------------
\section{Identificación de requisitos}
\label{sec:requisitos}
% -----------------------------------------------------------

El diseño de \textit{kubescan} parte de un análisis de requisitos derivado tanto de
las limitaciones identificadas en el estado del arte (Sección~\ref{sec:sintesis}) como
de los vectores de ataque documentados por la \ac{NSA} y la \ac{CISA} \citep{NSA2022}.

\subsection{Requisitos funcionales}

\begin{description}
  \item[RF-1 — Detección de misconfiguraciones individuales.] El sistema debe clasificar
    cada manifiesto \ac{YAML} como seguro o \textit{misconfigured}, produciendo además un
    índice de severidad en tres niveles (limpio, bajo\,--\,medio, alto\,--\,crítico).
  \item[RF-2 — Detección de cadenas de ataque.] El sistema debe clasificar un conjunto
    de manifiestos que describe un clúster en una de tres categorías: \textit{clean},
    \textit{isolated} (misconfiguraciones sin cadena viable) o
    \textit{attack\_chain} (ruta de escalada de privilegios hasta movimiento lateral).
  \item[RF-3 — Interfaz de línea de comandos.] El sistema debe exponer un comando
    \texttt{kubescan scan <dir>} para análisis de directorios locales y un comando
    \texttt{kubescan live} para análisis de clústeres activos mediante \textit{kubectl}.
  \item[RF-4 — Formatos de salida.] El sistema debe producir informes en formato texto
    legible y en formato \acs{JSON} estructurado para integración con herramientas externas.
  \item[RF-5 — Priorización de riesgos.] El sistema debe producir una clasificación
    ordenada de clústeres por puntuación de riesgo que maximice la Precisión@5, colocando las
    cadenas de ataque reales en las primeras posiciones.
\end{description}

\subsection{Requisitos no funcionales}

\begin{description}
  \item[RNF-1 — Trazabilidad.] Todas las reglas de detección de nodos de escape y
    movimiento lateral deben ser explícitas, auditables y derivadas de fuentes
    normativas documentadas \citep{NSA2022,BishopFox2021}.
  \item[RNF-2 — Reproducibilidad.] Los modelos entrenados deben almacenarse como
    \textit{checkpoints} e incluirse en el paquete, de forma que el análisis sea
    idempotente dado el mismo directorio de entrada.
  \item[RNF-3 — Instalabilidad.] El sistema debe ser instalable mediante
    \texttt{pip install -e kubescan/} sin dependencias de compiladores externos.
  \item[RNF-4 — Latencia.] El análisis de un clúster de hasta 500 manifiestos debe
    completarse en menos de 60 segundos en hardware de gama media (CPU).
\end{description}

% -----------------------------------------------------------
\section{Arquitectura general del sistema}
% -----------------------------------------------------------

El sistema \textit{kubescan} implementa una arquitectura de tres capas que transforma
un directorio de manifiestos \ac{YAML} en una puntuación de riesgo de cadena de ataque
(Figura~\ref{fig:arquitectura}):

\begin{enumerate}
  \item \textbf{Capa~1 (Random Forest):} extrae 28 características por manifiesto
    (25 indicadores binarios y 3 agregados derivados) y produce un
    \textit{risk\_score} \(\in [0,1]\).
  \item \textbf{Capa~2 (\ac{GAT}):} modela el clúster como un grafo donde los nodos
    son manifiestos y las aristas representan relaciones de seguridad; produce una
    probabilidad de cadena de ataque \(\in [0,1]\).
  \item \textbf{Capa~3 (Ensemble \acs{GA}):} combina las señales de las dos capas anteriores
    con una señal binaria de escape según la ecuación:
\end{enumerate}

\begin{equation}
  \text{score}(C) = w_{\text{rf}} \cdot \overline{\text{risk}}(C)
                  + w_{\text{gnn}} \cdot p_{\text{chain}}(C)
                  + w_{\text{esc}} \cdot \mathbf{1}[\exists\,\text{nodo de escape en }C]
  \label{eq:ensemble}
\end{equation}
\addequationtoindex{Función de puntuación del ensemble (Capa 3)}

donde los pesos \(w_{\text{rf}}, w_{\text{gnn}}, w_{\text{esc}}\) son optimizados
mediante Algoritmo Genético sobre las predicciones \textit{out-of-fold} de la
validación cruzada (Sección~\ref{sec:capa3_diseno}).

\begin{figure}[htbp]
  \centering
  \caption{Arquitectura del sistema \textit{kubescan}}
  \label{fig:arquitectura}
  \includegraphics[width=\textwidth]{fig_arquitectura}

  \fuentepropia
\end{figure}

La Figura~\ref{fig:arquitectura} muestra el proceso de tres capas, desde los
manifiestos \ac{YAML} de entrada hasta la clasificación de riesgo del clúster. Los
\textit{risk\_scores} producidos por la Capa~1 se incorporan como atributo nodal~26
del grafo de entrada de la Capa~2.

% -----------------------------------------------------------
\section{Construcción del conjunto de datos}
\label{sec:construccion_dataset}
% -----------------------------------------------------------

La construcción del conjunto de datos parte de la recolección de manifiestos
\ac{YAML} procedentes de fuentes públicas heterogéneas y avanza hasta la
extracción de las características que alimentan la Capa~1. Esta sección describe
primero el origen y la composición de los manifiestos recopilados y, a
continuación, el proceso de extracción de características que transforma cada
manifiesto en un vector estructurado.

\subsection{Fuentes de datos}
\label{sec:fuentes_datos}

El conjunto de datos tabular final contiene 2.827~filas y 58~columnas, combinando
cinco fuentes (agrupadas en cuatro categorías, dado que Rahman et~al.\ aporta dos
plataformas distintas con características de distribución diferentes; ver Tabla~\ref{tab:dataset}):

\begin{table}[htbp]
\centering\small
\caption{Composición del conjunto de datos de manifiestos Kubernetes}
\label{tab:dataset}
\begin{tabular}{@{}p{5.4cm}rp{4.9cm}@{}}
\toprule
Fuente & Filas & Rol \\
\midrule
Rahman et~al.\ \citep{Rahman2022} — GitHub & 1.510 & Manifiestos reales, etiquetados \\
Rahman et~al.\ \citep{Rahman2022} — GitLab & 396 & Ídem, plataforma distinta \\
BishopFox BadPods \citep{BishopFox2021} & 128 & Pods vulnerables intencionales \\
Kubernetes Goat \citep{KubernetesGoat} & 14 & Escenarios de vulnerabilidad \\
Repositorios de seguridad adicionales & 779 & CTFs, \textit{fixtures} de herramientas, labs \\
\midrule
\textbf{Total} & \textbf{2.827} & \\
\bottomrule
\end{tabular}
\fuentepropia
\end{table}

La incorporación de BadPods fue esencial para mitigar la escasez de ejemplos de
misconfiguraciones críticas en los repositorios reales: antes de añadir esta fuente,
\texttt{TRUE\_HOST\_PID} aparecía en sólo el 0,3\,\% de las filas; tras la incorporación,
la cobertura ascendió al 3,78\,\%. En cuanto a los conjuntos de datos de referencia
existentes para intrusión en red, como UNSW-NB15 \citep{Moustafa2015}, se estudió su
metodología de construcción como guía para el diseño del proceso de etiquetado y la
estrategia de balanceo de clases, dado que sus particiones curadas de entrenamiento y
evaluación presentan un desequilibrio de clases del mismo orden de magnitud que el de
este trabajo (56,3\,\% seguros vs.\ 43,7\,\% mal configurados). El conjunto \textit{GenKubeSec}
\citep{Malul2024}, publicado en 2024, no estaba disponible durante el desarrollo de
este trabajo; en su lugar se empleó \textit{Checkov} \citep{Checkov2021} directamente
sobre los ficheros \ac{YAML} descargados como fuente de validación multi-herramienta
equivalente.

\subsection{Características extraídas}

El extractor de características aplica 28 reglas de análisis estático a cada manifiesto
\ac{YAML}, organizadas en tres grupos:

\begin{itemize}
\raggedright
  \item \textbf{18 flags Rahman} (subconjunto del trabajo original): \texttt{TRUE\_HOST\_PID},
    \texttt{CAP\_SYS\_ADMIN}, \texttt{SEC\_CONT\_OVER\_PRIVIL}, \texttt{INSECURE\_HTTP},
    etc.
  \item \textbf{3 features derivadas}: \texttt{cap\_misuse} (OR de las dos flags de
    capacidades de sistema), \texttt{all\_secrets} (OR de los indicadores de exposición
    de secretos), \texttt{total\_misconfigs} (suma de todas las flags activas).
  \item \textbf{7 features extendidas}: \texttt{NO\_RUN\_AS\_NON\_ROOT},
    \texttt{IMAGE\_USES\_LATEST}, \texttt{SA\_AUTOMOUNT\_TOKEN}, \texttt{HOSTPATH\_MOUNT},
    etc.; imputadas mediante \textit{Checkov} para el 47,4\,\% de filas con YAML
    descargado, y mediante la mediana de columna para el resto.
\end{itemize}

La Capa~2 (\ac{GNN}) consume únicamente las 18~flags Rahman y las 7~features extendidas
(25 en total, \texttt{FEATURE\_COLS} en \texttt{yaml\_parser.py}) como vector de nodo,
excluyendo las 3~features derivadas —agregados calculados a partir de las otras 25 que
no aportan información estructural adicional al grafo—; ver
Sección~\ref{sec:features_table}.

Tres features presentan varianza casi nula pero se mantienen en el entrenamiento (con
importancia residual o nula, Sección~\ref{sec:rf_importance}):
\texttt{SECCOMP\_UNCONFINED} (15/2.827 activa), \texttt{VALID\_TAINT\_SECRET}
(0/2.827 activa) y \texttt{NO\_NETWORK\_POLICY} (2.818/2.827 activa, sin poder
discriminativo).

% -----------------------------------------------------------
\section{Construcción del grafo de clúster}
\label{sec:construccion_grafo}
% -----------------------------------------------------------

Cada directorio de manifiestos se representa, mediante \textit{NetworkX}
\citep{Hagberg2008}, como un grafo dirigido
\(G=(V, E)\) donde \(|V|\) es el número de manifiestos del clúster y
\(E\) contiene aristas de cinco tipos semánticamente distintos
(Tabla~\ref{tab:edge_types}). La Figura~\ref{fig:grafo_cluster} ilustra un ejemplo
con cuatro nodos que incluye los tres tipos de nodo posibles: escape, movimiento
lateral y limpio.

\begin{table}[htbp]
\centering\small
\caption{Tipos de aristas en el grafo de clúster}
\label{tab:edge_types}
\begin{tabular}{@{}p{3cm}cp{3.1cm}p{4.3cm}@{}}
\toprule
Tipo & Código & Dirección & Criterio \\
\midrule
Proximidad de directorio & 0 & Bidireccional & Mismo subdirectorio \ac{YAML} \\
Alcance de privilegio & 1 & Dirigida (escape\(\to\)todos) & Flag de escape activa \\
Movimiento lateral \acs{SA} & 2 & Dirigida (lateral\(\to\)todos) & Flag de movimiento lateral \\
Co-espacio de nombres (\texttt{SEMANTIC\_NS}) & 3 & Bidireccional & Mismo \textit{namespace} K8s \\
Escalada \ac{RBAC} & 4 & Dirigida (elevado\(\to\)todos) & \ac{SA} vinculado a rol privilegiado \\
\bottomrule
\end{tabular}
\fuentepropia
\end{table}

Cada nodo \(v_i\) lleva un vector de características \(\mathbf{x}_i \in \mathbb{R}^{26}\):
índices 0–24 corresponden a las 25 features binarias de la Capa~1 y el índice~25
contiene el \textit{risk\_score} producido por el modelo \textit{Random Forest}.

\begin{figure}[htbp]
  \centering
  \caption{Ejemplo de grafo de clúster Kubernetes con cuatro nodos}
  \label{fig:grafo_cluster}
  \includegraphics[width=0.72\textwidth]{fig_grafo_cluster}
  
  \fuentepropia
\end{figure}

La Figura~\ref{fig:grafo_cluster} muestra los tres tipos de nodo (escape, movimiento
lateral y limpio) y las aristas dirigidas que cada uno origina. El nodo rojo (pod-A) tiene
flags de escape activas y emite aristas de tipo~1 (\textit{privilege\_reach}) hacia
todos los demás nodos del clúster. El nodo naranja (pod-B) tiene flags de
movimiento lateral y emite aristas de tipo~2 (\textit{sa\_lateral}). La arista de
tipo~4 (\ac{RBAC}) representa la vinculación de un \textit{ServiceAccount} a un rol
privilegiado.

El conjunto de datos de grafos final comprende 599~grafos originales (25~clean, 537~isolated,
37~attack\_chain) y 1.154~grafos totales tras la aumentación de datos aplicada
exclusivamente en el conjunto de entrenamiento. Esta cifra incorpora el corpus
completo de Rahman et~al.\ \citep{Rahman2022}, cuyos \textit{namespaces} cruzados
se resolvieron mediante descarga autenticada desde GitHub, lo que amplió
sustancialmente el número de clústeres \textit{isolated} respecto a las fases
experimentales preliminares descritas en la Sección~\ref{sec:layer2} del
Capítulo~\ref{cap:resultados}.

\subsection{Detección de escape y movimiento lateral}

{\raggedright
La clasificación de nodos utiliza dos conjuntos de flags definidos en el código
como constantes inmutables (\texttt{frozenset}). Un nodo se clasifica como de
\textbf{escape} si tiene activa al menos una de las siguientes 8~flags:
\texttt{TRUE\_HOST\_PID}, \texttt{TRUE\_HOST\_IPC}, \texttt{TRUE\_HOST\_NET},
\texttt{DOCKERSOCK\_PATH}, \texttt{CAP\_SYS\_ADMIN}, \texttt{CAP\_SYS\_MODULE},
\texttt{SEC\_CONT\_OVER\_PRIVIL} o \texttt{HOSTPATH\_MOUNT}. Este último vector fue
incorporado durante el desarrollo al detectar que \textit{badpods\_hostpath} —que monta
el sistema de ficheros del anfitrión— era incorrectamente etiquetado como
\textit{isolated} al no tener ninguna otra flag de escape activa. Con la corrección,
el número de clústeres de cadena de ataque del corpus preliminar de 307~grafos
ascendió de 13 a 15 (25 en el corpus extendido de 471; 37 en el corpus final de 599).\par}

Un nodo se clasifica como de \textbf{movimiento lateral} si tiene activa al menos una
de las siguientes 4~flags: \texttt{SA\_AUTOMOUNT\_TOKEN} (token de cuenta de servicio
automontado), \texttt{USES\_DEFAULT\_SA} (usa la cuenta de servicio por defecto del
\textit{namespace}), \texttt{WITHIN\_MANIFEST\_SECRET} (secreto referenciado en el
manifiesto) o \texttt{ALLOW\_PRIVI} (contenedor con capacidad de escalada). Las aristas
de movimiento lateral (tipo~2) se construyen de forma análoga a las de escape: se emite
una arista dirigida desde cada nodo lateral hacia todos los demás, pero únicamente si
no existe ya una arista de mayor prioridad (escape o \ac{RBAC}).

\subsection{Detección de privilegio \acs{RBAC}}

La construcción de aristas \ac{RBAC} (tipo~4) requiere dos pasadas sobre el conjunto de
manifiestos. La primera pasada identifica los roles privilegiados: aquellos cuyo nombre
está en el conjunto \{cluster-admin, admin, edit\} o cuyas reglas contienen verbos de
escalada \{*, escalate, bind, impersonate\}. La segunda pasada recorre los objetos
\textit{RoleBinding} y \textit{ClusterRoleBinding}: sólo cuando el campo
\textit{roleRef} referencia un rol privilegiado identificado en la primera pasada se
emite una arista \ac{RBAC} desde el pod que usa esa cuenta de servicio hacia todos los demás
nodos del clúster.

% -----------------------------------------------------------
\section{Capa 1: Random Forest}
\label{sec:capa1_diseno}
% -----------------------------------------------------------

La Capa~1 clasifica cada manifiesto de forma individual a partir de las
características descritas en la sección anterior. Se entrenan dos modelos
independientes sobre el mismo conjunto de características: un clasificador
binario que distingue manifiestos seguros de mal configurados, y un modelo de
severidad que estima el grado de gravedad de la configuración incorrecta
detectada. A continuación se describe el diseño de cada uno.

\subsection{Modelo binario}

Se entrena un \textit{RandomForestClassifier} \citep{Breiman2001} de \textit{scikit-learn}
\citep{Pedregosa2011} con los hiperparámetros siguientes:

\begin{itemize}
  \item \texttt{n\_estimators = 500}
  \item \texttt{max\_features = sqrt} (clasificación estándar)
  \item \texttt{min\_samples\_leaf = 2} (evita hojas de muestra única)
  \item \texttt{class\_weight = balanced} (compensa ratio 1,29:1)
  \item \texttt{random\_state = 42}
\end{itemize}

El modelo se entrena sobre 2.261~muestras (80\,\% del total) y se evalúa en 566
muestras de test no vistas durante el entrenamiento. La probabilidad predicha para la
clase misconfigured (clase~1) se usa como \textit{risk\_score} del nodo en el grafo.

\subsection{Modelo de severidad}

Además del modelo binario, se entrena un clasificador de severidad de tres clases
(\textit{clean}, \textit{low-medium}, \textit{high-critical}) para ofrecer información
más granular al operador de seguridad. Este modelo no interviene en el cálculo del
grafo; se incluye como salida adicional de la Capa~1.

% -----------------------------------------------------------
\section{Capa 2: Graph Attention Network}
\label{sec:capa2_diseno}
% -----------------------------------------------------------

La Capa~2 opera sobre el grafo de clúster construido en la sección anterior y
clasifica el clúster completo según el riesgo de cadena de ataque que presenta.
Esta sección describe la arquitectura del modelo \acs{GAT}, el procedimiento de
entrenamiento empleado y la estrategia de aumentación de datos aplicada para
mitigar la escasez de clústeres con cadena de ataque confirmada.

\subsection{Arquitectura}

El modelo \ac{GAT} \citep{Velickovic2018} implementado tiene la siguiente estructura,
con las dimensiones exactas verificadas en \texttt{gat\_encoder.py}:

\begin{enumerate}
\raggedright
  \item \textbf{Embedding de tipo de arista}: \texttt{Embedding(5, 8)}.
        Los cinco tipos de arista (proximidad de directorio, alcance de
        privilegio, lateral por \textit{ServiceAccount}, espacio de nombres y
        \ac{RBAC}) son categóricos: codificarlos como un escalar continuo
        implicaría un orden espurio (el tipo~4 «valdría» cuatro veces el
        tipo~1). Cada tipo se proyecta a un vector aprendido de 8~dims que
        \texttt{GATConv} consume como atributo de arista
        (\textit{edge\_dim=8}).
  \item \textbf{Proyección inicial}: \texttt{Linear(26, 256)} \(+\) \texttt{LayerNorm(256)}
        \(+\) \acs{ReLU} \(+\) \texttt{Dropout(0.3)}.
        Proyecta el vector de nodo de 26~dims a 256~dims
        (\(64 \times 4\,\text{cabezas}\)).
  \item \textbf{Capas \ac{GAT}~0 y~1} (\textit{concat=True}):
        \texttt{GATConv(256, 64, heads=4)} con \textit{edge\_dim=8};
        salida 256~dims \((64 \times 4)\); seguidas de \texttt{LayerNorm}, \acs{ELU},
        conexión residual y \texttt{Dropout(0.3)}.
  \item \textbf{Capa \ac{GAT}~2} (\textit{concat=False, heads=1}):
        \texttt{GATConv(256, 64, heads=1)}; salida 64~dims
        (una sola cabeza de atención que produce directamente
        el embedding final del nodo, sin concatenación);
        seguida de \texttt{LayerNorm}, \ac{ELU} y \texttt{Dropout(0.3)}
        (sin conexión residual: las dimensiones de entrada y salida difieren,
        256 frente a 64).
  \item \textbf{Global pooling}: \texttt{global\_mean\_pool} y \texttt{global\_max\_pool}
        sobre los 64-dim embeddings de todos los nodos; concatenación
        \(\to\) 128~dims (\(64 + 64\)).
  \item \textbf{Clasificador \ac{MLP}}: \texttt{Linear(128, 64)} \(+\) \ac{ReLU} \(+\)
        \texttt{Dropout(0.3)} \(+\) \texttt{Linear(64, 3)}.
\end{enumerate}

La elección de \ac{GAT} sobre \ac{GCN} responde a que el mecanismo de atención permite al modelo
aprender qué aristas son más informativas para la señal de ataque: una arista de
\textit{privilege\_reach} debe recibir mayor peso que una arista de proximidad de
directorio. El \textit{pooling} combinado de media y máximo captura dos señales
complementarias: la media refleja la postura de seguridad promedio del clúster,
mientras que el máximo identifica el nodo más peligroso, ambas necesarias para detectar
cadenas de ataque que pueden estar concentradas en un único nodo crítico.

\subsection{Entrenamiento}

El modelo se entrena con \textit{PyTorch Geometric} \citep{Fey2019}, construido
sobre \textit{PyTorch} \citep{Paszke2019}, bajo los siguientes hiperparámetros: optimizador AdamW con tasa de aprendizaje
\(5\times10^{-4}\) y \textit{weight\_decay}\,=\,\(10^{-4}\);
\textit{scheduler} CosineAnnealingLR con T\_max\,=\,300 y
\(\eta_{\min}\,=\,10^{-5}\); función de pérdida CrossEntropyLoss con pesos
de clase calculados por frecuencia inversa;
\textit{gradient clipping} con norma máxima 2,0; \textit{early stopping} con
paciencia de 50 épocas; tamaño de lote 8; semilla global 42 (idéntica para el
entrenamiento y para la generación de particiones); y validación cruzada
estratificada de 5~pliegues. El \textit{checkpoint} de cada pliegue se selecciona
por la mejor Precisión@5 de validación, con la F1 macro como criterio de
desempate (\texttt{--select-by p5}), alineando la selección de modelos con la
métrica primaria del sistema en lugar de con la métrica de clasificación.

Los pesos de clase se calculan en cada pliegue de entrenamiento mediante la
siguiente fórmula implementada en \texttt{train\_gnn.py}:

\begin{equation}
  w_i = \frac{1/n_i}{\displaystyle\sum_{j=0}^{K-1} 1/n_j} \cdot K
  \label{eq:class_weights}
\end{equation}
\addequationtoindex{Cálculo de pesos de clase para el entrenamiento de la GAT}

donde \(n_i\) es el número de grafos de entrenamiento de la clase \(i\) y \(K=3\)
es el número de clases. Esta fórmula normaliza los pesos para que su media sea~1,0,
proporcionalmente equivalente al comportamiento de \texttt{class\_weight=balanced} de
\textit{scikit-learn} (misma ponderación relativa entre clases). Cabe destacar que los pesos se calculan
\textbf{después} de la aumentación de datos: como la aumentación genera 15~variantes
por grafo de cadena de ataque, la clase \textit{attack\_chain} pasa de ser la
minoría a ser la mayoría en el conjunto de entrenamiento (aproximadamente 380~grafos
aumentados vs.\ 18~clean y 367~isolated en el conjunto de entrenamiento de cada
pliegue, tras excluir las variantes de los clústeres de validación y test). La Ecuación~\ref{eq:class_weights} se auto-adapta a esta nueva distribución
asignando mayor peso a \textit{clean} e \textit{isolated}, y menor peso a
\textit{attack\_chain}, compensando el efecto de la sobre-representación creada por
la aumentación.

\subsection{Aumentación de datos}

Para mitigar la escasez de ejemplos de cadena de ataque (37 de 599 grafos
originales), se aplican tres estrategias de aumentación exclusivamente sobre los
grafos de entrenamiento: \textit{edge dropout} (eliminación aleatoria de entre el
10\,\% y el 35\,\% de aristas de los tipos~0 y~3, los menos críticos para la señal
de ataque, con tasa muestreada uniformemente por variante);
\textit{feature masking} (puesta a cero, con probabilidad por nodo entre el 10\,\%
y el 25\,\%, restringida a seis características ajenas a la regla de etiquetado de
cadena — nunca las flags de escape, \texttt{HOSTPATH\_MOUNT} ni el \textit{risk\_score}); y
\textit{subgraph sampling} (extracción de un subgrafo inducido por entre el 40\,\%
y el 70\,\% de nodos, aplicada a las últimas cinco de las 15~variantes y solo a
grafos de más de 50~nodos, con \textit{BFS}
sembrada en el nodo de mayor \textit{risk\_score}; cada subgrafo
se valida contra la regla de etiquetado de cadena y, si deja de cumplirla, se
descarta en favor de una variante sobre el grafo completo, evitando introducir
ruido de etiqueta). Cada grafo de
cadena de ataque origina 15~variantes aumentadas, resultando en 555~grafos
adicionales (37~originales \(\times\) 15~variantes) en el \textit{pool} de
aumentación; de éstos, únicamente 390 (los correspondientes a los 26~grafos de
cadena cuya base pertenece a la partición de entrenamiento) entran efectivamente
en el entrenamiento, para un total de 1.154~grafos con aumentación.

La asignación de variantes a particiones es \textbf{consciente de grupos}
(\textit{group-aware}): una variante aumentada solo puede entrar en una partición
de entrenamiento si su grafo base pertenece a esa misma partición. Las variantes
derivadas de clústeres de validación o de test se excluyen del entrenamiento por
completo, ya que entrenar con casi-duplicados de un grafo de evaluación
constituiría fuga de datos.

El particionado es además \textbf{consciente de familias de plantilla}: los
clústeres originales que comparten una misma plantilla base (por ejemplo, las
siete variantes \texttt{badpods\_*}, que difieren únicamente en qué flag de
escape activan, o las nueve \texttt{datadog\_*}) son casi duplicados entre sí, no
muestras independientes. Cuando una familia contiene etiquetas mixtas, sus
miembros se asignan como una unidad atómica a una única partición — nunca se
reparten entre entrenamiento, validación y test ni entre pliegues — y la
asignación emplea un reparto voraz ponderado por el número de grafos de cada
etiqueta, de modo que la proporción real de cada clase por partición siga las
fracciones solicitadas. Sin esta regla, un modelo entrenado con cinco hermanos de
una familia resuelve trivialmente al sexto en test, inflando las métricas. La
partición resultante (semilla 42) es de 815~grafos de entrenamiento (con
aumentación), 88 de validación y 86 de test. Adicionalmente, los 86~clústeres de
test se mantienen fuera de los pliegues de validación cruzada: los pliegues
particionan únicamente los 513~grafos originales restantes (32 de cadena), de
modo que ni los modelos de pliegue ni las predicciones \textit{out-of-fold}
empleadas para ajustar los pesos del \textit{ensemble} (Capa~3) observan ningún
clúster de test, directa ni indirectamente.

% -----------------------------------------------------------
\section{Capa 3: Ensemble con Algoritmo Genético}
\label{sec:capa3_diseno}
% -----------------------------------------------------------

Los pesos \((w_{\text{rf}}, w_{\text{gnn}}, w_{\text{esc}})\) de la
Ecuación~\ref{eq:ensemble} son optimizados mediante un Algoritmo Genético
\citep{Holland1975} con 60~individuos codificados como pares
\((w_{\text{gnn}}, w_{\text{esc}})\) \(\in [0,1]^2\), con \(w_{\text{rf}}\) derivado
como \(1 - w_{\text{gnn}} - w_{\text{esc}}\) (la representación 2D garantiza que la
suma de los tres pesos sea siempre~1). La función objetivo es:

\begin{equation}
  F = 0{,}7 \cdot P@5 + 0{,}3 \cdot (1 - \text{FPR\_clean})
      - \lambda \sum_{i} \max(0,\ w_{\min} - w_i)
  \label{eq:ga_fitness}
\end{equation}
\addequationtoindex{Función objetivo (fitness) del Algoritmo Genético}

El último término es una penalización de suelo con \(w_{\min}\,=\,0{,}1\) y
\(\lambda\,=\,2{,}0\): sin ella, el optimizador es libre de anular cualquiera de
los tres pesos cuando eso puntúa marginalmente mejor sobre la muestra
\textit{out-of-fold}, lo que produce soluciones degeneradas que no generalizan al
conjunto de test. La penalización hace que ignorar una señal por completo esté
estrictamente dominado salvo que la ganancia observada supere \(\lambda\) veces el
déficit.

La selección se realiza mediante torneo de tamaño~3, con elitismo del 10\,\%; el cruce
es aritmético (media de los dos padres) con probabilidad 0,7; la mutación combina
perturbación gaussiana (\(\sigma\,=\,0{,}08\), tasa del 15\,\%) con una mutación de
frontera (tasa del 5\,\%) que salta a los vértices del símplex, seguida de
re-proyección al dominio válido (\(w_{\text{rf}} \geq 0\)); la población inicial se
siembra con los vértices y puntos medios del símplex y con muestreo de Dirichlet, y el
criterio de parada es 150~generaciones.
Como \textit{baseline} determinístico, se ejecuta previamente una búsqueda en rejilla
2D con paso \(1/20\) (231~evaluaciones); la comparación de su óptimo con el resultado
del \ac{GA} se discute en la Sección~\ref{sec:layer3}.

La señal de escape se define como un indicador binario que vale~1,0 si al menos un
nodo del clúster tiene alguna flag de escape activa, y~0,0 en caso contrario. Esta
formulación evita la dilución que produciría una fracción cuando la mayoría de los
nodos del clúster son limpios.

La asignación de clase final a partir de la puntuación \textit{ensemble\_score}
se realiza mediante los siguientes umbrales fijos, definidos en
\texttt{ga\_ensemble.py}:
\begin{itemize}
  \item \textit{ensemble\_score}\,\(\geq\)\,0,60 \(\ \Rightarrow\) \texttt{ATTACK\_CHAIN}
  \item \textit{ensemble\_score}\,\(\geq\)\,0,30 \(\ \Rightarrow\) \texttt{ISOLATED\_MISCONFIG}
  \item De lo contrario \(\Rightarrow\) \texttt{CLEAN}
\end{itemize}
Estos umbrales son fijos y no se reoptimizan durante el entrenamiento del \ac{GA},
que optimiza únicamente los pesos \(w_i\) de la Ecuación~\ref{eq:ensemble}.
Permiten al operador interpretar directamente la puntuación: cualquier clúster con
\textit{score}\,\(\geq\)\,0,60 requiere revisión inmediata.

\subsection{Restricción de viabilidad estructural}
\label{sec:gate_estructural}

Una cadena de ataque multi-salto requiere, por definición, alcanzabilidad entre al
menos dos manifiestos; un grafo de clúster con un único nodo no puede albergarla.
Esta propiedad, verificada sobre el corpus completo (los 592~grafos de cadena
—37 originales y 555 variantes aumentadas— tienen \(\geq\)2 nodos y \(\geq\)1 arista,
mientras que 499 de los 537 grafos
\textit{isolated} son de un solo nodo), se incorpora como restricción de
inferencia en \texttt{ga\_ensemble.py} (\texttt{MIN\_CHAIN\_NODES},
\texttt{is\_chain\_feasible}, \texttt{chain\_rank\_key}): en el
\textit{ranking} por riesgo de cadena, ningún clúster de un solo manifiesto puede
superar a un clúster estructuralmente viable, y el veredicto de un clúster
inviable se limita a \texttt{ISOLATED\_MISCONFIG} aunque su puntuación supere el
umbral de 0,60. La restricción se aplica únicamente en inferencia y evaluación;
el ajuste de los pesos del \ac{GA} se realiza sobre el \textit{ranking} sin
restricción, donde los pesos están mejor determinados (la Sección~\ref{sec:layer3}
cuantifica el efecto de esta regla sobre las métricas de test).

% -----------------------------------------------------------
\section{Implementación de kubescan}
\label{sec:implementacion}
% -----------------------------------------------------------

El sistema se distribuye como un paquete Python instalable mediante
\texttt{pip install -e kubescan/}. La estructura del paquete sigue la convención
\textit{src-layout} de \textit{setuptools}: el código fuente reside en
\texttt{kubescan/src/kubescan/} y los módulos se organizan en tres submódulos:
\texttt{model/} (clasificadores \ac{RF}, \ac{GAT} y ensamblador \ac{GA}), \texttt{utils/}
(\textit{parser} \ac{YAML} y constructor de grafos) y \texttt{cli.py} (interfaz de
línea de comandos basada en \textit{Click}).

El punto de entrada principal expone dos comandos. El primero,
\texttt{kubescan scan <dir>}, analiza un directorio local de manifiestos \ac{YAML} y
emite un informe de riesgo en formato texto o \ac{JSON}. El segundo,
\texttt{kubescan live}, consulta el \ac{API} de \textit{Kubernetes} mediante
\textit{kubectl}, vuelca los recursos de \textit{workloads} y \ac{RBAC} en un directorio
temporal y ejecuta el \textit{pipeline} completo sobre ellos. La resolución de los
pesos del modelo sigue la cadena: argumento \texttt{--checkpoints-dir}
\(\to\) variable de entorno \texttt{KUBESCAN\_CHECKPOINTS} \(\to\) enlace simbólico
\texttt{kubescan/checkpoints/trained/} \(\to\) \texttt{FileNotFoundError}.

% -----------------------------------------------------------
\section{Tabla de características del sistema}
\label{sec:features_table}
% -----------------------------------------------------------

Las Tablas~\ref{tab:features_a} y~\ref{tab:features_b} describen las 28 características
utilizadas por la Capa~1, organizadas por grupo y con indicación de si son relevantes
para la detección de escape (ESC) o de movimiento lateral (LAT) en el grafo
de la Capa~2. De estas 28, las 25 que no son derivadas (18~Rahman + 7~extendidas)
también forman el vector de nodo de la Capa~2.

\begin{table}[htbp]
\centering\small
\caption{Características del clasificador \ac{RF} — grupo Rahman (1/2)}
\label{tab:features_a}
\begin{tabular}{@{}p{4.2cm}llp{4.5cm}@{}}
\toprule
Nombre & Grupo & Rol en grafo & Descripción \\
\midrule
\texttt{TRUE\_HOST\_PID} & Rahman & ESC & Comparte PID namespace del anfitrión \\
\texttt{TRUE\_HOST\_IPC} & Rahman & ESC & Comparte IPC namespace del anfitrión \\
\texttt{TRUE\_HOST\_NET} & Rahman & ESC & Comparte red del anfitrión \\
\texttt{DOCKERSOCK\_PATH} & Rahman & ESC & Monta socket de Docker \\
\texttt{CAP\_SYS\_ADMIN} & Rahman & ESC & Capacidad \texttt{SYS\_ADMIN} activa \\
\texttt{CAP\_SYS\_MODULE} & Rahman & ESC & Capacidad \texttt{SYS\_MODULE} activa \\
\texttt{WITHIN\_MANIFEST\_SECRET} & Rahman & LAT & Secreto en texto plano embebido en el manifiesto \\
\texttt{SEC\_CONT\_OVER\_PRIVIL} & Rahman & ESC & Contenedor privilegiado (\texttt{privileged: true}) \\
\texttt{ALLOW\_PRIVI} & Rahman & LAT & \texttt{allowPrivilegeEscalation: true} \\
\texttt{INSECURE\_HTTP} & Rahman & — & Usa HTTP sin TLS en puerto conocido \\
\texttt{NO\_SECU\_CONTEXT} & Rahman & — & Ausencia de \texttt{securityContext} definido \\
\texttt{HOST\_ALIAS} & Rahman & — & Usa \texttt{hostAliases} para resolución DNS manual \\
\texttt{NO\_DEFAULT\_NSPACE} & Rahman & — & Despliega en el \textit{namespace} \texttt{default} \\
\texttt{NO\_RESO} & Rahman & — & Sin límites/\textit{requests} de recursos definidos \\
\texttt{NO\_ROLLING\_UPDATE} & Rahman & — & Sin estrategia de actualización \textit{RollingUpdate} \\
\texttt{SECCOMP\_UNCONFINED} & Rahman & — & Perfil \textit{seccomp} \texttt{unconfined} (sin filtrado de \textit{syscalls}) \\
\texttt{VALID\_TAINT\_SECRET} & Rahman & — & Secreto válido referenciado en \textit{taint}/\textit{toleration} \\
\texttt{NO\_NETWORK\_POLICY} & Rahman & — & Sin \texttt{NetworkPolicy} definida en el \textit{namespace} \\
\bottomrule
\end{tabular}
\fuentepropia
\end{table}

\begin{table}[htbp]
\centering\small
\caption{Características del clasificador \ac{RF} — grupos derivado y extendido, y nodo (2/2)}
\label{tab:features_b}
\begin{tabular}{@{}p{4.2cm}llp{4.5cm}@{}}
\toprule
Nombre & Grupo & Rol en grafo & Descripción \\
\midrule
\texttt{cap\_misuse} & Derivada & — & OR lógico de \texttt{CAP\_SYS\_ADMIN} y \texttt{CAP\_SYS\_MODULE} \\
\texttt{all\_secrets} & Derivada & — & OR de los indicadores de exposición de secretos \\
\texttt{total\_misconfigs} & Derivada & — & Suma de todas las flags activas \\
\texttt{NO\_RUN\_AS\_NON\_ROOT} & Extendida & — & Sin restricción de usuario no root \\
\texttt{NO\_READ\_ONLY\_ROOT\_FS} & Extendida & — & Sistema de ficheros raíz con escritura \\
\texttt{IMAGE\_USES\_LATEST} & Extendida & — & Imagen con \texttt{tag: latest} \\
\texttt{SA\_AUTOMOUNT\_TOKEN} & Extendida & LAT & \textit{ServiceAccount} con token automontado \\
\texttt{USES\_DEFAULT\_SA} & Extendida & LAT & Usa la \textit{ServiceAccount} \texttt{default} del \textit{namespace} \\
\texttt{UNTRUSTED\_REGISTRY} & Extendida & — & Imagen proveniente de un registro no confiable \\
\texttt{HOSTPATH\_MOUNT} & Extendida & ESC & Monta ruta del sistema de ficheros del anfitrión \\
\texttt{risk\_score} & Nodo (Capa~2) & — & Probabilidad \ac{RF} clase misconfigured \(\in [0,1]\) \\
\bottomrule
\end{tabular}
\fuentepropia
\end{table}

Las 8~flags de escape (ESC, \texttt{ESCAPE\_FLAGS} en \texttt{graph\_builder.py}) son
\texttt{TRUE\_HOST\_PID}, \texttt{TRUE\_HOST\_IPC}, \texttt{TRUE\_HOST\_NET},
\texttt{DOCKERSOCK\_PATH}, \texttt{CAP\_SYS\_ADMIN}, \texttt{CAP\_SYS\_MODULE},
\texttt{SEC\_CONT\_OVER\_PRIVIL} y \texttt{HOSTPATH\_MOUNT}. Las 4~flags de movimiento
lateral (LAT, \texttt{LATERAL\_FLAGS}) son \texttt{SA\_AUTOMOUNT\_TOKEN},
\texttt{USES\_DEFAULT\_SA}, \texttt{WITHIN\_MANIFEST\_SECRET} y \texttt{ALLOW\_PRIVI}.
Las 7~características extendidas están imputadas mediante \textit{Checkov} para el
47,4\,\% de filas con YAML disponible y mediante la mediana de columna para el resto.

% -----------------------------------------------------------
\section{Calidad de software y proceso de validación}
\label{sec:calidad}
% -----------------------------------------------------------

El desarrollo de \textit{kubescan} sigue prácticas de ingeniería de software
orientadas a la producción. La calidad del código se garantiza mediante cuatro
herramientas integradas en un \textit{pipeline} de pre-commit y \acs{CICD}: (1)
\textit{ruff}, linter y formateador de código Python con reglas de estilo y
seguridad; (2) \textit{mypy} en modo estricto, para comprobación estática de
tipos en todos los módulos del paquete; (3) \textit{pytest}, para la suite de
tests unitarios e de integración del paquete; y (4) \textit{pre-commit hooks},
que aplican automáticamente ruff y mypy en cada commit.

{\raggedright
La suite de tests cubre los siguientes componentes: el parser de \ac{YAML}
(\texttt{utils/yaml\_parser.py}), la construcción del grafo de clúster
(\texttt{utils/graph\_builder.py}), el clasificador \ac{RF} (\texttt{model/rf\_classifier.py})
y el ensamblador \ac{GA} (\texttt{model/ga\_ensemble.py}). Los tests de integración
verifican el \textit{pipeline} completo de extremo a extremo: dado un directorio
de manifiestos de prueba con flags de escape conocidas, se comprueba que la salida
del comando \texttt{kubescan scan} produce el veredicto esperado (ATTACK\_CHAIN)
con \textit{ensemble\_score} por encima del umbral de clasificación.\par}
